\chapter{Pre-test and Post-test Actions, and Equipment}\label{ch:actions}

A schedule can do work before the first step and after the last one. These are Pre-test and
Post-test Actions: not steps, so they have no Step Number, no result and no inputs.

\begin{itemize}
  \item \texttt{:setPreTestAction(fn)} runs once before the first Test Step. It is the place to
    open a Device, declare Equipment and read the bench configuration.
  \item \texttt{:setPostTestAction(fn)} runs once after the last step, in the reverse of the
    Pre-test order, and is called with the run's overall result: \texttt{"PASS"} or
    \texttt{"FAIL"}.
  \item A Pre-test Action that raises aborts the run before any step runs. The run records
    \texttt{PRETEST\_FAILED} and every step \texttt{NOTRUN}.
  \item Post-test Actions run even after an abort. ATE clears the abort flag before them so a
    run can always tear down; an abort that lands during one Action ends that Action and leaves
    the rest to run. An Action that errors fails the run and the others still run.
\end{itemize}

Both Actions are called with the run's standing result as their first argument. For a Pre-test
Action that is always \texttt{"PASS"}; nothing has had a chance to fail yet.

\noindent\textbf{\texttt{actions.tslua}}\par
\lstinputlisting[language=ATELua]{../examples/test-writer/actions.tslua}

\section{Equipment}

\texttt{ATE.addEquipment(name, serial, connection, calDue, notes)} records a piece of test
equipment for the Test Report's Equipment section. Every argument is a string, only
\texttt{name} is required, and \texttt{calDue} is free text that ATE prints as given without
checking it. Declaring the same name and serial twice adds no second row. Equipment is not a
Device: ATE never links the two, and a schedule may list Equipment it never connects to.

\section{Logging, and reading the run's own context}

\texttt{ATE.log(level, \dots)} writes a line into the Test Run's log; the arguments are turned
to text and joined with spaces. \texttt{level} is one of \texttt{Log.DEBUG},
\texttt{Log.INFO}, \texttt{Log.WARNING} or \texttt{Log.ERROR}. The main screen shows INFO and
above; every level is stored. \texttt{print} is redirected to \texttt{ATE.log(Log.INFO, \dots)},
so a stray print is still recorded.

\texttt{ATE.run} is a read-only table describing the run in progress. Its fields are:

\begingroup\footnotesize
\begin{verbatim}
UUT_serial, Operator, Station, Project, Test_Schedule_Name,
Media_Release_Media_Number, name, version
\end{verbatim}
\endgroup

Writing to any of them is an error.

\noindent\textbf{\texttt{logging.tslua}}\par
\lstinputlisting[language=ATELua]{../examples/test-writer/logging.tslua}

\cando{open a Device and declare the bench in a Pre-test Action, log what happened, and close
things up in a Post-test Action that runs whatever the run did.}
